% page 561 du fac-simile (image p-560 du scan)
% bloc 172.7 x 237.7 mm ; marge gauche 29.3 mm
\pgc{7.620mm}{0.000mm}{
\l{\cel{54}553.}
\l{}
\l{\sou{sterleto}. (\sou{zool.)}\cel{2}Speco di sturgo, quan onu renkontras en la}
\l{\cel{3}parto orientala di Europa ed en Azia, e qua aspektas kom la}
\l{\cel{3}formo yuna di la sturgo. De olu onu obtenas kaviaro. - L.}
\l{\cel{3}\sou{acipensar ruthenus}. - DEFR.}
\l{}
\l{\sou{sternar}.\cel{2}(\sou{trans.}) Extensar (persono, kozo, su) longesale sur}
\l{\cel{3}ulo. - L.}
\l{}
\l{\sou{sternumo}.\cel{2}(\sou{anat.})\cel{2}Osto plata an-avan torako, meze di olca, e}
\l{\cel{3}sur qua \textquotedbl{}muntesas\textquotedbl{} la kosti. - DEFIS.}
\l{}
\l{\sou{sternutar}. (\sou{netrans.})\cel{2}Expirar bruske, tra nazo e boko, per movo}
\l{\cel{3}konvulsatra di diafragmo. - EFIS.}
\l{}
\l{\sou{stertorar}. (\sou{netrans.})\cel{2}Agar ta bruiso ne-normala quan produktas,}
\l{\cel{3}en la respir-organi, la paso di aero tra mukosi quin sekrecis la}
\l{\cel{3}bronki, la pulmoni. - SL.}
\l{}
\l{\sou{stetoskopo}. Sorto di tubo akustikala, quan onu uzas por la auskulti}
\l{\cel{3}medikala. - DEFIS.}
\l{}
\l{\sou{stifa}. Inter-parenta pro la mariajo dato-duesma di un ek la ge-}
\l{\cel{3}patri. - D.}
\l{}
\l{\sou{stifto}. (\sou{tekn.}) Kelo qua asemblas la peci ye qui konsistas char-}
\l{\cel{3}niro, buklo, e c. - DR.}
\l{}
\l{\sou{stigmato}. I. marko neefacebla quan lasis plago, bruluro, vunduro.II.}
\l{\cel{3}- (r\sou{eligio kristana}).Marko ek la kin vunduri di la korpo di Iesu}
\l{\cel{3}Kristo, imprimita mirakle sur la korpo di ula santi. - III.}
\l{\cel{3}Marko imprimita, per fera stampilo redigita per varmigo, sur}
\l{\cel{3}la shultro-pelo di ula krimininti. - IV. (\sou{bot.}) Korpo glandatra}
\l{\cel{3}ube finas la stilo di pistilo, e transmisas a lu, por igar olu}
\l{\cel{3}pasar aden la ovario, la parto fekundigiva di poleno.- DEFIRS.}
\l{}
\l{\sou{stilo}. I. (\sou{epoki antiqua}). Puncilo di qua un extremajo pintoza}
\l{\cel{3}uzesis por trasar la literi sur vaxo-plako, e di qua la altra,}
\l{\cel{3}platigite, uzesis por efacar. - II. (\sou{literaturo)}\cel{2}Maniero ex-}
\l{\cel{3}presar redakte sua pensi, sua idei. - III. (\sou{arto}). Karaktero}
\l{\cel{3}generala di la verki da artisto, skolo, yarcento, naciono. - IV.}
\l{\cel{3}\sou{(bot.)}\cel{2}Prolonguro filoforma di la somito di la ovario, quan}
\l{\cel{3}finas la stigmato. - DEFRS.}
\l{}
\l{\sou{stileto}. Poniardo, kun lamo streta, secione triangula o quadri-}
\l{\cel{3}angula. - EF.}
\l{}
\l{\sou{stilistiko.}\cel{2}Regularo pri stilo, qua esas situebla inter\cel{2}gramatlko\textquotedbl{}}
\l{\cel{3}e \textquotedbl{}retoriko\textquotedbl{}.- (\sou{segun Bally}). Havante ideo klara o sentimento,}
\l{\cel{3}la arto expresar li per la moyeni quin disponas linguo e quaze}
\l{\cel{3}inter-sinonima.-DF.}
\l{}
\l{\sou{stilobato}.(\sou{arkitekt.}) Basamento piedestalo-forma qua suportas ko-}
\l{\cel{3}loni. - DEFIS.}
}
